%%%%%%%%%%%%%%%%%%%%%%%%%%%%%%%%%%%%%%%%%
% Medium Length Professional CV
% LaTeX Template
% Version 2.0 (8/5/13)
%
% This template has been downloaded from:
% http://www.LaTeXTemplates.com
%
% Original author:
% Rishi Shah
%
% Important note:
% This template requires the resume.cls file to be in the same directory as the
% .tex file. The resume.cls file provides the resume style used for structuring the
% document.
%
%%%%%%%%%%%%%%%%%%%%%%%%%%%%%%%%%%%%%%%%%

%----------------------------------------------------------------------------------------
%	PACKAGES AND OTHER DOCUMENT CONFIGURATIONS
%----------------------------------------------------------------------------------------

\documentclass{resume} % Use the custom resume.cls style

\usepackage[left=0.75in,top=0.6in,right=0.75in,bottom=0.6in]{geometry} % Document margins
\newcommand{\tab}[1]{\hspace{.2667\textwidth}\rlap{#1}}
\newcommand{\itab}[1]{\hspace{0em}\rlap{#1}}
\usepackage{xcolor}
\definecolor{cvlink}{HTML}{173B63}
\usepackage[colorlinks=true,linkcolor=cvlink,urlcolor=cvlink,citecolor=cvlink]{hyperref}
\name{Fran\c{c}ois Bares} % Your name
\address{Sa\"{i}d Business School, University of Oxford \\ Park End Street, Oxford OX1 1HP, United Kingdom} % Your address
\address{\href{mailto:francois.bares.dphil@said.oxford.edu}{email} \\ (+44) 7759 335022 \\ \href{https://fbares.github.io}{website}} % Your email, phone number and website

\begin{document}

%----------------------------------------------------------------------------------------
%	EDUCATION SECTION
%----------------------------------------------------------------------------------------

\begin{rSection}{Education}

% If the degree is not yet awarded when you send this, write "2022--2026 (expected)".
{\bf DPhil (PhD) in Finance, Sa\"{i}d Business School, University of Oxford} \hfill {\em 2022--2026} \\
Thesis: \textit{Three Essays on the Production of Knowledge} \\
Committee: \href{https://www.sbs.ox.ac.uk/about-us/people/thomas-hellmann}{Thomas Hellmann} (chair), \href{https://www.hbs.edu/faculty/Pages/profile.aspx?facId=951639}{Kyle Myers}, \href{https://www.sbs.ox.ac.uk/about-us/people/seung-joo-lee}{Seung Joo Lee} \\
Green Templeton College; full scholarship \\
{\bf Visiting PhD Student, Harvard Business School} \hfill {\em 2025} \\
Hosted by Kyle Myers; affiliated with the Laboratory for Innovation Science at Harvard (LISH) \\
{\bf MS in Applied Economics, University of Wisconsin--Madison} \hfill {\em 2020--2022} \\
First two years of the PhD programme; full scholarship \\
{\bf MSc in Economics, London School of Economics} \hfill {\em 2017--2018}

\end{rSection}

%----------------------------------------------------------------------------------------
%	RESEARCH
%----------------------------------------------------------------------------------------


% \begin{rSection}{Ongoing Research}

% {\bf The allocative efficiency of research funding} \\
% \textit{Solo-authored} \\
% How should research funding should be allocated across research fields? I explore this question in the context of medical research in the U.S. Using data from the National Institute of Health (NIH) I consider the allocative efficiency of NIH funding across disease. The prevailing view in the literature advocates funding disease proportionally to their social cost, implicitly assuming equal marginal return of public funding across diseases. I test this assumption by estimating the marginal product for each disease, taking into account knowledge spillovers and a "researcher stealing effects" whereby different diseases compete for a fixed supply of research labor. I do so with a structural model of research production with disease choices \textit{a la} Roy, augmented with disease-level knowledge externalities. Early estimates suggest that 1) the marginal product of funding varies greatly across diseases, 2) Cross-disease knowledge spillovers are slightly larger and exhibit a far greater variance than own disease spillovers and 3) Positive knowledge externalities outweigh negative "researchers stealing" effect for most disease, suggesting that omitting externalities underestimate returns on investment.

% {\bf Large funding shocks in science} \\
% \textit{Solo-authored} \\
%  Public efforts to boost innovation in a specific sector, sometimes called "mission-oriented", "big push", or "moonshot" policies, have grown popular in recent years, often as part of broader industrial policies. This paper develops a microfounded framework to evaluate the impact of such policies on sectoral production in science. A famous example is the 1971 National Cancer Act, better known as the "war on cancer" that quadrupled U.S. cancer research funding over ten years. Using the methodology pioneered by David Baqaee and Emmanuel Farhi, I model how funding shocks propagate through a network of knowledge exchange. In this framework, large shocks differ from small shocks because they are large enough for the envelope theorem to break down, researchers to re-optimise and change the knowledge they produce and consume.


% %{\bf AI-qualising Science} \\
% %\textit{With Anantha Divakaruni and Ludovic Phalippou} \\
% %An artificial intelligence (AI) system called AlphaFold has revolutionised protein research by accurately predicting three-dimensional protein structures, earning its developers the 2024 Nobel Prize in Chemistry. Researchers in lower-ranked universities adopt AlphaFold more. Post-AlphaFold, they publish more in top journals relative to researchers in top universities. The results are driven by increased collaboration between researchers from high-ranked and low-ranked universities. This study highlights the role of collaboration networks in cementing inequality in science, and how a technological shock can reshape these networks.

% {\bf Technological uncertainty and innovation finance, evidence from Xprize competitions} \\
% \textit{With Thommas Hellmann} \\
% XPrize is the world’s largest innovation prize organiser. Xprize competitions lead to an exogenous reduction in the uncertainty surrounding specific firms and broader innovation sectors. Leveraging XPrize’s unique setting and data, this study assesses the role that technological uncertainty plays on the funding of early-stage start-ups. Using a combination of event studies and regression discontinuity designs, we estimate the effect of technological uncertainty for both public and private funders, at the sectoral and firm level.

% \end{rSection}


\begin{rSection}{Job Market Paper}

\parbox[t]{0.82\linewidth}{\textbf{\href{https://fbares.github.io/assets/papers/where-should-science-funding-go.pdf}{``Where Should Science Funding Go? Estimating Marginal Returns to Public Funding Across Research Fields''}}} \hfill {\em 2026} \\
Science funding is shrinking in many countries, while commercial innovation increasingly relies on publicly funded science. I study the optimal allocation of science funding across fields. I 
estimate marginal returns to funding in 19 biomedical fields using a micro-founded structural model where scientists compete for grants. These scientists differ in productivity and shift their research focus towards better-funded fields. Once funded, they produce knowledge by combining grant funding with existing knowledge from their own and other fields, giving rise to knowledge spillovers. Marginal returns vary widely across fields: an additional \$1 million yields 2.2 publications in Nursing Research and 14.8 in Human Genome. Knowledge spillovers account for most of this gap. Upstream fields, such as General Medicine and Human Genome, generate far larger spillovers than specialised downstream fields, such as Nursing Research or Deafness and Communication Disorders. Scientists do move towards better-funded fields, but their replacements are equally productive once funded, so this reallocation costs little in equilibrium. I construct three welfare-relevant measures of scientific output: academic impact, commercial potential and disease-burden reduction. Marginal-return rankings barely change across the three measures. I combine each measure with the estimated marginal returns from the model to compare the current allocation with two counterfactual allocations: allocating funding in proportion to disease burden, as the medical community often advocates, and in proportion to marginal returns, as the model suggests. Across the three measures, allocation based on the model estimates yields 39--54\% more output than the current allocation and 15.1--26.4\% more than disease-burden alignment.

\textit{Presented at:} Oxford FAME Seminar; NBER Productivity Seminar; BDI--EIEF--CEPR Conference on Finance and Innovation.

\end{rSection}


\begin{rSection}{Other Papers}

\parbox[t]{0.82\linewidth}{\textbf{\href{https://fbares.github.io/assets/papers/innovation-prizes-as-fundraising-signals.pdf}{``Innovation Prizes as Fundraising Signals: Evidence from XPRIZE Competitions''}}} \hfill {\em 2026} \\
With Thomas Hellmann, Elaine Hungenberg, and Andrew Tauhert

\item \parbox[t]{0.82\linewidth}{\textbf{\href{https://fbares.github.io/assets/papers/does-ai-democratize-science.pdf}{``Does Artificial Intelligence Democratize Science? Evidence from AlphaFold''}}} \hfill {\em 2026} \\
With Anantha Divakaruni and Ludovic Phalippou. Previously circulated as ``AI-qualizing Science''

\item \parbox[t]{0.82\linewidth}{\textbf{\href{https://link.springer.com/article/10.1057/s41267-026-00860-z}{``The Impact of the Global Minimum Tax on Incentives for Business Location, Investment, and Profit Shifting''}}} \hfill {\em 2026} \\
With Michael Devereux, Irem Guceri, and Vikram Patil \\
\textit{Journal of International Business Studies}, 57(7), 972--989

\item \parbox[t]{0.82\linewidth}{\textbf{\href{https://www.sbs.ox.ac.uk/sites/default/files/2020-02/OECD_GloBE_proposal_report_2020.pdf}{``The OECD Global Anti-Base Erosion Proposal''}}} \hfill {\em 2020} \\
With Michael Devereux, Sarah Clifford, Judith Freedman, Irem Guceri, Martin McCarthy, Martin Simmler, and John Vella \\
\textit{Report for PwC}

\end{rSection}


\begin{rSection}{Work in Progress}

\parbox[t]{0.82\linewidth}{\textbf{``The Cost of Hormones: RCT Gender Balance and Drug Development''}}

\item \parbox[t]{0.82\linewidth}{\textbf{``The Social Returns to R\&D: Evidence from Cancer Research''}} \hfill \\
With Kyle Myers and Mai Anh Tran

\item \parbox[t]{0.82\linewidth}{\textbf{``Capital Intensity and Risk Sharing in Deep-Tech Startups''}} \hfill \\
With Kyle Myers and Mai Anh Tran

\end{rSection}


% Uncomment and fill in: conference and seminar presentations are a standard section on a job-market CV.
% \begin{rSection}{Presentations}
%
% {\em 2026:} FAME seminar (Oxford), ... \\
% {\em 2025:} ...
%
% \end{rSection}


\begin{rSection}{Teaching Experience}

{\bf Sa\"{i}d Business School, University of Oxford} \hfill {\em 2024--2026} \\
	Applied Financial Economics: \textit{Support class lecturer} (master's, 2026) \\
	Executive Finance: \textit{Support class lecturer} (MBA, 2024, 2026) \\
	Financial Accounting: \textit{Support class lecturer} (MBA, 2025) \\
	Financial Strategy: \textit{Teaching assistant} (executive education, 2025) \\
{\bf Stanford University in Oxford} \hfill {\em 2024} \\
	Game Theory and Financial Economics: \textit{Tutor} (undergraduate) \\
{\bf Oriel College, University of Oxford} \hfill {\em 2023, 2024} \\
	Economics of Information: \textit{Summer school lecturer} (undergraduate, 2023, 2024) \\
	Corporate Finance: \textit{Summer school lecturer} (undergraduate, 2024) \\
{\bf University of Wisconsin--Madison} \hfill {\em 2022} \\
	World Hunger and Malnutrition: \textit{Teaching assistant} (undergraduate)

\end{rSection}

%----------------------------------------------------------------------------------------
%	WORK EXPERIENCE SECTION
%----------------------------------------------------------------------------------------





\begin{rSection}{Other Research Experience}

\textbf{Research Assistantships} \\
Ramana Nanda, Imperial College London \hfill {\em 2022--2023} \\
Thomas Hellmann, University of Oxford \hfill {\em 2022--2023} \\
Andrew Stevens, University of Wisconsin--Madison \hfill {\em 2021} \\
Paul Frijters, London School of Economics \hfill {\em 2017--2018}

{\bf Centre for Business Taxation, University of Oxford} \hfill {\em 2018--2020} \\
\textit{Pre-doctoral Fellow} \\
Worked on the tax treatment of corporate losses in the UK with large-scale administrative data. \\
Worked on business rates and their effects on commercial property use. Paper available \href{https://papers.ssrn.com/sol3/papers.cfm?abstract_id=4352363}{here}. \\
% Co-authored a \href{https://www.sbs.ox.ac.uk/sites/default/files/2020-02/OECD_GloBE_proposal_report_2020.pdf}{report} quantifying the implications of the OECD's pillar two tax reform proposal. \\
% Co-authored a study on the distortionary effect of a global minimum tax rate.
% \\
Co-authored a study on the optimality of combining a patent box with an R\&D tax credit policy.


% {\bf 89 Initiative Think Tank, London School of Economics} \hfill {\em 2016-2017} \\
% \textit{Policy Officer} \\
% Background research on relevant academic and legal literature. \\
% Drafted policy proposals that were then presented to the European Parliament. \\

{\bf Cr\'{e}dit Agricole S.A., Economic Research Division} \hfill {\em 2016} \\
\textit{Summer Intern} \\
Specified and estimated macro-level models of corporate investment for the Eurozone and six of its member states (VAR models) and built the underlying database.

\end{rSection}



\begin{rSection}{Fellowships and Awards}

Research Fellowship in Finance, Sa\"{i}d Business School \hfill {\em 2026} \\
Graduate Scholarship, Sa\"{i}d Business School \hfill {\em 2022} \\
Green Templeton College--Sa\"{i}d Business School DPhil Scholarship \hfill {\em 2022} \\
Graduate Fellowship, University of Wisconsin--Madison \hfill {\em 2020}

\end{rSection}



\begin{rSection}{Languages and Technical Background}

\begin{tabular}{ @{} >{\bfseries}l @{\hspace{3ex}} l }
Programming \ & R, Stata, Python, LaTeX, Unix, Docker, Git, Query DSL, SQL \\
Data \ & Orbis, PATSTAT, Eurostat, OpenAlex, Dimensions.ai, NIH RePORTER, \\ \ & Google Patents, PitchBook, Nielsen Retail Scanner \\
 \ & Accredited by the UK government to use confidential administrative data \\
Modelling \ & Game theory, discrete choice, supply chain networks, production functions \\
Statistics \ & Econometrics (parametric and non-parametric), machine learning (NLP) \\
Languages \ & French (native), English (fluent, C2), German (B2), Spanish (A2) \\
\end{tabular}

\end{rSection}


% Uncomment once your letter writers are confirmed; add their email addresses.
% \begin{rSection}{References}
%
% {\bf Thomas Hellmann} (chair), Sa\"{i}d Business School, University of Oxford \\
% {\bf Kyle Myers}, Harvard Business School \\
% {\bf Seung Joo Lee}, Sa\"{i}d Business School, University of Oxford
%
% \end{rSection}


\begin{rSection}{Recommendations}

\textbf{Thomas Hellmann} \\
\textit{DP World Professor in Entrepreneurship and Innovation} \\
Sa\"{i}d Business School, University of Oxford, Park End Street, Oxford OX1 1HP, UK \\
\href{mailto:thomas.hellmann@sbs.ox.ac.uk}{thomas.hellmann@sbs.ox.ac.uk}

\smallskip
\textbf{Kyle Myers} \\
\textit{Anna Spangler Nelson and Thomas C. Nelson Associate Professor of Business Administration} \\
Harvard Business School, Boston, MA 02163, USA \\
\href{mailto:kmyers@hbs.edu}{kmyers@hbs.edu}


\smallskip
\textbf{Seung Joo Lee} \\
\textit{Associate Professor of Finance} \\
Sa\"{i}d Business School, University of Oxford, Park End Street, Oxford OX1 1HP, UK \\
\href{mailto:seung.lee@sbs.ox.ac.uk}{seung.lee@sbs.ox.ac.uk}

\end{rSection}

\end{document}
